% Section 5.1: Definition by relations | Part: 1 | Depends on: sec:metric, sec:ext | Status: done | Original lines: 1286-1316
%!TEX root = ../../main.tex
\subsection{Definition by relations}

\begin{definition}[Clifford algebra]
Let $\mathcal I_Q\subset T(V)$ be the two-sided ideal generated by $\{v\otimes v-Q(v)\,1: v\in V\}$. The
\emph{Clifford algebra} is $\Cl(V,Q)=T(V)/\mathcal I_Q$. Its product is called the
\emph{geometric product} and is written by juxtaposition: $uv$. Thus
\begin{equation}\label{eq:cliff}
 v\,v=Q(v)=v\cdot v\qquad\text{for every }v\in V.
\end{equation}
When $V=\R^{p,q}$ we write $\Cl(p,q)$; the algebra is then also called the \emph{geometric algebra} $\GA(V)$ and its elements \emph{multivectors}.
\end{definition}

\begin{theorem}[Universal property]\label{thm:cliffuniv}
Let $A$ be a unital associative algebra and $f:V\to A$ linear with $f(v)^2=Q(v)1_A$ for all $v$. Then $f$ extends uniquely to an algebra homomorphism $\Cl(V,Q)\to A$.
\end{theorem}
\begin{proof}
By Theorem~\ref{thm:Tuniv}, $f$ extends to $\bar f:T(V)\to A$; the hypothesis says $\bar f$ kills each generator $v\otimes v-Q(v)$ of $\mathcal I_Q$, hence all of $\mathcal I_Q$, so $\bar f$ descends to the quotient.
\end{proof}

\begin{proposition}[Fundamental identity]\label{prop:fund}
For $u,v\in V$: $\;uv+vu=2\,g(u,v)=2\,u\cdot v$. In particular, orthogonal vectors anticommute.
\end{proposition}
\begin{proof}
Polarize \eqref{eq:cliff}: $(u+v)^2-u^2-v^2=uv+vu$, and by \eqref{eq:polar} this equals $2g(u,v)$.
\end{proof}

Hence, for an orthonormal basis $(e_i)$ with $e_i^2=\varepsilon_i=\pm1$, the \emph{entire} algebra is generated by the relations
\[
 e_ie_j=-e_je_i\ (i\neq j),\qquad e_i^2=\varepsilon_i .
\]
